\chapter{The UpstreamDrift Platform}
\label{ch:platform}

\epigraph{%
  The best theory is one you can run.
}{}

\section{Architecture Overview}

UpstreamDrift is a unified platform for multi-body dynamics simulation, trajectory
optimization, and motor control analysis. It wraps multiple physics engines behind
a consistent Python API.

\subsection{Core Components}

\begin{center}
\renewcommand{\arraystretch}{1.3}
\begin{tabular}{@{}lll@{}}
\toprule
\textbf{Component} & \textbf{Location} & \textbf{Purpose} \\
\midrule
Physics Engines & \texttt{src/engines/physics\_engines/} & Multi-body simulation \\
Shared Infrastructure & \texttt{src/engines/common/} & State, export, diagnostics \\
Pendulum Models & \texttt{src/engines/pendulum\_models/} & Simplified test cases \\
Learning Modules & \texttt{src/learning/} & ML integration \\
Reinforcement Learning & \texttt{src/reinforcement\_learning/} & RL benchmarks \\
Robotics & \texttt{src/robotics/} & Robot control modules \\
API & \texttt{src/api/} & External access layer \\
Tools & \texttt{src/tools/} & Development utilities \\
\bottomrule
\end{tabular}
\end{center}

\subsection{Supported Physics Engines}

\begin{enumerate}
  \item \textbf{MuJoCo} (\texttt{physics\_engines/mujoco/}): Google DeepMind's
        multi-joint dynamics engine. Soft contacts, muscle actuators, fast simulation.
        Primary engine for biomechanical modeling.

  \item \textbf{Drake} (\texttt{physics\_engines/drake/}): MIT's multibody dynamics
        toolbox. Mathematical rigor, analytical gradients, trajectory optimization.
        Primary engine for control design.

  \item \textbf{Pinocchio} (\texttt{physics\_engines/pinocchio/}): LAAS-CNRS C++ library
        with Python bindings. Efficient rigid-body dynamics with analytical derivatives.
        Primary engine for real-time control.

  \item \textbf{OpenSim} (\texttt{physics\_engines/opensim/}): Stanford's musculoskeletal
        modeling software. Hill-type muscle models, inverse dynamics. Primary engine
        for biomechanics.

  \item \textbf{MyoSuite} (\texttt{physics\_engines/myosuite/}): Muscle-driven
        reinforcement learning environments. Primary engine for motor learning research.
\end{enumerate}

\section{Getting Started}

\begin{lstlisting}[language=Python, caption={Basic UpstreamDrift usage pattern.}]
"""UpstreamDrift: Getting started with multi-engine simulation.

This script demonstrates the basic workflow:
1. Load a model
2. Configure simulation parameters
3. Run a simulation
4. Extract and analyze results
"""
import numpy as np
from pathlib import Path

# Configuration
MODEL_PATH = Path("src/engines/pendulum_models")
ENGINE = "mujoco"  # or "drake", "pinocchio"

def create_simple_pendulum(
    length: float = 1.0,
    mass: float = 1.0,
    damping: float = 0.1
) -> dict:
    """Create a simple pendulum model configuration.

    Parameters
    ----------
    length : float
        Pendulum length in meters.
    mass : float
        Bob mass in kilograms.
    damping : float
        Joint damping coefficient.

    Returns
    -------
    dict
        Model configuration dictionary.
    """
    return {
        "name": "simple_pendulum",
        "n_dof": 1,
        "segments": [{
            "name": "bob",
            "mass": mass,
            "length": length,
            "com_position": length / 2,
            "inertia": mass * length**2 / 3,
        }],
        "joints": [{
            "name": "pivot",
            "type": "revolute",
            "axis": [0, 0, 1],
            "damping": damping,
        }],
        "gravity": [0, -9.81, 0],
    }


def simulate_pendulum(
    config: dict,
    q0: float = np.pi / 4,
    t_end: float = 5.0,
    dt: float = 0.001
) -> dict:
    """Simulate a simple pendulum using Euler integration.

    Parameters
    ----------
    config : dict
        Model configuration.
    q0 : float
        Initial angle (radians from vertical).
    t_end : float
        Simulation duration (seconds).
    dt : float
        Time step.

    Returns
    -------
    dict
        Simulation results with keys 't', 'q', 'qd', 'energy'.
    """
    seg = config["segments"][0]
    m = seg["mass"]
    L = seg["length"]
    b = config["joints"][0]["damping"]
    g = abs(config["gravity"][1])
    I = seg["inertia"]

    n_steps = int(t_end / dt)
    t = np.zeros(n_steps)
    q = np.zeros(n_steps)
    qd = np.zeros(n_steps)
    energy = np.zeros(n_steps)

    q[0] = q0
    qd[0] = 0.0

    for k in range(n_steps - 1):
        # Equations of motion: I * qdd + b * qd + m*g*L/2*sin(q) = 0
        qdd = (-b * qd[k] - m * g * L / 2 * np.sin(q[k])) / I

        # Semi-implicit Euler
        qd[k + 1] = qd[k] + qdd * dt
        q[k + 1] = q[k] + qd[k + 1] * dt
        t[k + 1] = t[k] + dt

        # Energy
        KE = 0.5 * I * qd[k + 1]**2
        PE = m * g * L / 2 * (1 - np.cos(q[k + 1]))
        energy[k + 1] = KE + PE

    return {"t": t, "q": q, "qd": qd, "energy": energy}


# Run simulation
config = create_simple_pendulum(length=1.0, mass=1.0, damping=0.05)
results = simulate_pendulum(config, q0=np.pi/3, t_end=10.0)

print(f"Simulation complete: {len(results['t'])} timesteps")
print(f"Initial energy: {results['energy'][0]:.4f} J")
print(f"Final energy:   {results['energy'][-1]:.4f} J")
print(f"Energy loss:    {(1 - results['energy'][-1]/results['energy'][0])*100:.1f}%")
\end{lstlisting}

\section{Engine Selection Guide}

\begin{handson}{Choosing the Right Engine}{engine-choice}
  \begin{center}
  \renewcommand{\arraystretch}{1.3}
  \begin{tabular}{@{}lll@{}}
  \toprule
  \textbf{If you need\ldots} & \textbf{Use} & \textbf{Why} \\
  \midrule
  Fast contact simulation & MuJoCo & Soft contact model, GPU \\
  Trajectory optimization & Drake & Analytical gradients \\
  Real-time forward dynamics & Pinocchio & C++ speed, O(n) \\
  Muscle-driven biomechanics & OpenSim & Hill models built-in \\
  RL with muscles & MyoSuite & Gym-compatible \\
  Quick prototyping & Built-in Python & No dependencies \\
  \bottomrule
  \end{tabular}
  \end{center}
\end{handson}

\section*{Exercises}

\begin{enumerate}
  \item \textbf{First Simulation.} Run the pendulum simulation and plot angle vs.\ time
        and energy vs.\ time. Verify that energy decreases due to damping.

  \item \textbf{Integration Method.} Replace the semi-implicit Euler with RK4. Compare
        energy conservation over 100 seconds.

  \item \textbf{Double Pendulum.} Extend the code to a double pendulum. Plot the
        chaotic trajectory in configuration space.

  \item \textbf{Engine Comparison.} Run the same pendulum simulation in MuJoCo and
        in the built-in Python code. Compare results quantitatively.
\end{enumerate}
